% Part: normal-modal-logic
% Chapter: tableaux
% Section: simple-S5

\documentclass[../../../include/open-logic-section]{subfiles}

\begin{document}

\olfileid{nml}{tab}{s5}

\olsection{\Log{S5} కోసం సరళ \usetoken{P}{tableau}}

\Log{S5}, సార్వత్రిక నమూనాల వర్గానికి సంబంధించి
నిర్దుష్టమైనదీ సంపూర్ణమైనదీ. సార్వత్రిక నమూనాలో ప్రతి
లోకం నుంచి ప్రతి లోకం ప్రాప్యమే. అటువంటి నమూనాల్లో
ప్రాప్యత సంబంధాన్ని ప్రత్యేకంగా గుర్తుంచుకోవాల్సిన పని
లేదు: ``$\mSat{M}{!A}[w]$ అయ్యే లోకం~$w$ ఉంది'' అనేది,
ఆ లక్షణం గల~$u$ నుంచి ప్రాప్యమైన లోకం~$w$ ఉంది అనేదానికి
సమానార్థకం. కనుక \Log{S5}లో నమూనాను లోకాల సమితి,
మూల్యనిర్దేశం~$V$ ద్వారానే నిర్వచించవచ్చు. దాంతో
\tetoken{టాబ్లో}{!!{tableau}} నియమాలనూ సరళీకరించవచ్చని
తెలుస్తుంది. సాధారణ సందర్భంలో పూర్వసూచికలుగా ధన
పూర్ణసంఖ్యల అనుక్రమాలను తీసుకుంటాం; అప్పుడు ఏ
పూర్వసూచిక ఏ పూర్వసూచిక పేరుగల లోకం నుంచి ప్రాప్యమైన
లోకాన్ని సూచిస్తుందో గుర్తుంచుకోవచ్చు: $\sigma.n$ అనేది
$\sigma$ పేరుగల లోకం నుంచి ప్రాప్యమైన లోకానికి పేరు.
కానీ \Log{S5}లో ప్రతి లోకం నుంచి ప్రతి లోకం ప్రాప్యం;
కాబట్టి ఆ వివరాన్ని గుర్తుంచుకోవాల్సిన అవసరం లేదు.
దానికి బదులు ధన పూర్ణసంఖ్యలనే పూర్వసూచికలుగా వాడవచ్చు.
సరళీకృత నియమాలు \olref{tab:rules-S5}లో ఉన్నాయి.

\begin{table}
  \begin{center}
    \def\arraystretch{3}\def\fCenter{}
    \begin{tabular}{|c|c|}
    \hline
    \iftag{prvBox}{
      \AxiomC{\sFmla{\True}{\Box !A}[n]}
      \RightLabel{\TRule{\True}{\Box}}
      \UnaryInfC{\sFmla{\True}{!A}[m]}
      \DisplayProof
      &
      \AxiomC{\sFmla{\False}{\Box !A}[n]}
      \RightLabel{\TRule{\False}{\Box}}
      \UnaryInfC{\sFmla{\False}{!A}[m]}
      \DisplayProof\\
      $m$ ఉపయోగించినది & $m$ కొత్తది
      \\[1ex]
      \hline}{}
    \iftag{prvDiamond}{
      \AxiomC{\sFmla{\True}{\Diamond !A}[n]}
      \RightLabel{\TRule{\True}{\Diamond}}
      \UnaryInfC{\sFmla{\True}{!A}[m]}
      \DisplayProof
      &
      \AxiomC{\sFmla{\False}{\Diamond !A}[n]}
      \RightLabel{\TRule{\False}{\Diamond}}
      \UnaryInfC{\sFmla{\False}{!A}[m]}
      \DisplayProof\\
      $m$ కొత్తది & $m$ ఉపయోగించినది
      \\[1ex]
      \hline}{}
    \end{tabular}
  \end{center}
  \caption{\Log{S5} కోసం సరళీకృత నియమాలు.}
  \ollabel{tab:rules-S5}
\end{table}

\begin{ex}
  $\Log{S5} \Proves \Ax{5}$, అంటే
  $\Diamond!A \lif \Box\Diamond!A$ అని చూపే సరళీకృత
  సంవృత టాబ్లోను ఇస్తున్నాం.
  \begin{oltableau}
    [\pFmla{\False}{\Diamond\formula{A} \lif \Box\Diamond \formula{A}}{1},
      just = \TAss
      [\pFmla{\True}{\Diamond \formula{A}}{1}, just = {\TRule{\False}{\lif}[1]}
        [\pFmla{\False}{\Box\Diamond \formula{A}}{1},
          just = {\TRule{\False}{\lif}[1]}
          [\pFmla{\False}{\Diamond \formula{A}}{2},
            just = {\TRule{\False}{\Box}[3]}
            [\pFmla{\True}{\formula{A}}{3},
              just = {\TRule{\True}{\Diamond}[2]}
                [\pFmla{\False}{\formula{A}}{3},
                  just = {\TRule{\False}{\Diamond}[4]}, close]
            ]
          ]
        ]
      ]
    ]
  \end{oltableau}
\end{ex}

\end{document}
